\section{Training budgets determined by economic primitives}
\label{sec:r27primitive}
The finite construction in Theorem~\ref{thm:r25finitepolicy} uses a bound on each
own-policy continuation once that continuation has been evaluated. We now
bound every such target before the first fit. This distinction matters for a
method-level statement: the input to the guarantee is an economy and a welfare
tolerance, not a favorable terminal training event.

Consider the additive capital economy of Section~\ref{sec:r23policy}.
Let $J^{\rm ref}>0$ bound the zero-control policy's value on
$\|x_0\|^2\leq b$, and set $\eta=\epsilon_0/J^{\rm ref}$. Define, using only
primitives,
\begin{equation}\label{eq:r27envelope}
 \overline P_T=Q_f,\qquad
 \overline P_t=(1+\eta)Q_t+\beta A_t'\overline P_{t+1}A_t.
\end{equation}
For $j<T$, choose $q_jI\preceq Q_j$, $r_tI\preceq R_t$, and
$L_j\geq d\|\overline P_j\|_2$, and write $m_j=dq_j$.
For a decision at $t=j-1$, let $a_t\geq\|A_t\|_2$,
$b_t\geq\|B_t\|_2$, and put
\begin{equation}\label{eq:r27threshold}
 \chi_t=\beta b_ta_t\left(1+\frac{\beta b_t^2L_j}{dr_t}\right),
 \qquad e_j=\frac{d\sqrt{\eta q_tr_t}}{\chi_t}.
\end{equation}
When $\chi_t=0$, the corresponding action is independent of the continuation;
no learning accuracy is required for that decision. The terminal continuation
is known exactly.

\begin{theorem}[A primitive-only finite NBO schedule]\label{thm:r27primitive}
Suppose the reference value is finite and the stage coercivity bounds are
positive. At each nonterminal continuation date $j$, initialize
$W_{j,0}=\sqrt{c_j}O_j$, where $O_j'O_j=I$ and $0<c_j\leq m_j$, and use
$\alpha_j=1/(4L_j)$. All hidden entries are updated by gradient descent on
$\frac14\|W'W-dP_j\|_F^2$, where $P_j$ evaluates the already finalized future
policy. It is sufficient to allow
\begin{equation}\label{eq:r27cap}
 N_j=\left\lceil
 \frac{\log^+(\sqrt d\,L_j/e_j)}{-\log(1-2\alpha_jc_j)}
 \right\rceil
\end{equation}
updates, then construct the current action as in
Theorem~\ref{thm:r25finitepolicy}. The returned policy satisfies
$J^\pi(x_0)-J^*(x_0)\leq\epsilon_0$ throughout $\|x_0\|^2\leq b$,
against every adapted finite-cost policy. Every target obeys
$m_jI\preceq dP_j\preceq d\overline P_j\preceq L_jI$.
For square factors and dense algebra, a sufficient real-arithmetic work account
is $O(d^3[T+\sum_{j=1}^{T-1}N_j])$, with $O(Td^2)$ stored matrix entries.
\end{theorem}

The proof is a backward induction, not a uniform-convergence assertion for an
unspecified optimizer. A local defect at most $\eta x'Q_tx$ implies
$P_t\preceq(1+\eta)Q_t+\beta A_t'P_{t+1}A_t$ by testing zero current control.
This closes the envelope induction while the isotropic contraction closes the
training induction. Both inductions refer to the same implemented future.
The full proof is in Section~\ref{sec:r27primitiveproof}.

The work statement counts real arithmetic. It is not a bit-complexity theorem
and does not remove finite-precision verification, memory traffic, or the cost
of evaluating a different economy. Conventional quadratic dynamic programming
can exploit the same primitives without training a factor. Thus this theorem
establishes a deterministic finite-work property of a specified trainable NBO,
not a lower bound on competing methods. It replaces a condition on a successful
training run by a condition on the economic input, without turning the
observed conventional advantage into a neural one.
